% ============================================================
% TUTORIAL: Fase 7 — Deploy final + verificacion end-to-end
% Proyecto: Sentinel AcademIA
% Compilar: tectonic tutorial-fase7-deploy.tex
% ============================================================

\documentclass[11pt,a4paper,oneside]{article}

\usepackage[utf8]{inputenc}
\usepackage[T1]{fontenc}
\usepackage[spanish,es-nodecimaldot]{babel}
\usepackage{lmodern}
\usepackage[margin=2.5cm]{geometry}
\usepackage{parskip}
\usepackage{microtype}

\usepackage{xcolor}
\usepackage{colortbl}
\usepackage{array}
\usepackage{booktabs}
\usepackage{amssymb}
\usepackage{pifont}
\usepackage{tcolorbox}
\tcbuselibrary{breakable,skins}

\usepackage{listings}

\lstdefinelanguage{yaml}{
  basicstyle=\ttfamily\small,
  keywordstyle=\color{blue!70!black}\bfseries,
  stringstyle=\color{red!70!black},
  commentstyle=\color{green!50!black}\itshape,
  morekeywords={Type,Properties,Version,Statement,Effect,Principal,Action,Resource},
}

\lstdefinelanguage{bash}{
  basicstyle=\ttfamily\small,
  keywordstyle=\color{blue!70!black}\bfseries,
  stringstyle=\color{red!70!black},
  commentstyle=\color{green!50!black}\itshape,
  morekeywords={aws,sam,npm,cd,ls,mkdir,rm,cp,mv,export,unset,echo,cat,curl,sleep,npx,git,python3,pip},
}

\lstdefinelanguage{json}{
  basicstyle=\ttfamily\small,
  showstringspaces=false,
  breaklines=true,
  literate=
    *{\{}{{{\{}}}{1}
     {\}}{{{\}}}}{1}
     {[}{{{[}}}{1}
     {]}{{{]}}}{1}
     {:}{{{:}}}{1}
     {,}{{{,}}}{1},
  string=[s]{"}{"},
  stringstyle=\color{red!70!black},
  morestring=[b]",
}

\lstset{
  basicstyle=\ttfamily\small,
  numberstyle=\tiny\color{gray},
  numbers=left,
  numbersep=8pt,
  frame=single,
  framerule=0.4pt,
  rulecolor=\color{gray!50},
  backgroundcolor=\color{gray!5},
  breaklines=true,
  breakatwhitespace=true,
  showstringspaces=false,
  captionpos=b,
  tabsize=2,
  literate={á}{{\'a}}1 {é}{{\'e}}1 {í}{{\'i}}1 {ó}{{\'o}}1 {ú}{{\'u}}1
           {ñ}{{\~n}}1 {Á}{{\'A}}1 {É}{{\'E}}1 {Í}{{\'I}}1 {Ó}{{\'O}}1 {Ú}{{\'U}}1
           {¿}{{?`}}1 {¡}{{!`}}1
}

\usepackage[colorlinks=true,linkcolor=blue!60!black,urlcolor=blue!60!black,citecolor=blue!60!black]{hyperref}

\usepackage{fancyhdr}
\pagestyle{fancy}
\fancyhf{}
\fancyhead[L]{\small\textsc{Sentinel AcademIA}}
\fancyhead[R]{\small Tutorial Fase 7}
\fancyfoot[C]{\thepage}
\renewcommand{\headrulewidth}{0.4pt}

\definecolor{okgreen}{RGB}{40,140,60}
\definecolor{errred}{RGB}{200,50,50}
\definecolor{warnorange}{RGB}{220,130,30}
\definecolor{infoblue}{RGB}{30,90,180}

\newcommand{\ok}{\textcolor{okgreen}{\textbf{\checkmark}}}
\newcommand{\err}{\textcolor{errred}{\textbf{\ding{55}}}}
\newcommand{\warn}{\textcolor{warnorange}{\textbf{\textasciitilde}}}

\title{\Huge\textbf{Fase 7: Deploy Final} \\[0.3em]
       \Large Verificaci\'on end-to-end + README \\[0.5em]
       \normalsize Tutorial paso a paso}
\author{Proyecto Sentinel AcademIA \\ \small lFunknown}
\date{\today}

\begin{document}

\maketitle
\thispagestyle{empty}

\begin{tcolorbox}[colback=blue!5!white,colframe=infoblue,title={\textbf{¿Qué hicimos en este tutorial?}}]
En la \textbf{Fase 7} (la \'ultima) hicimos el deploy final y verificamos que \textbf{todas las features funcionan end-to-end}. Este tutorial documenta:

\begin{itemize}
  \item \textbf{Re-deploy} con todos los fixes de bugs encontrados en testing
  \item \textbf{Test E2E completo}: 5 endpoints probados en orden
  \item \textbf{Listado de todos los recursos} del stack final
  \item \textbf{README.md} completo para el repositorio
  \item \textbf{Estado final}: 50/50 tests pasando, todos los servicios activos
\end{itemize}

\textbf{Hallazgos del deploy}:
\begin{itemize}
  \item \ok \textbf{7 Lambdas} desplegadas y funcionando
  \item \ok \textbf{DynamoDB} con tabla multi-tenant
  \item \ok \textbf{SQS + DLQ} para procesamiento as\'incrono
  \item \ok \textbf{S3} con bucket para evidencia + frontend
  \item \ok \textbf{API Gateway} con CORS configurado
  \item \ok \textbf{Frontend Vue 3} accesible p\'ublicamente
  \item \ok \textbf{OCI Object Storage} con knowledge base
\end{itemize}
\end{tcolorbox}

\begin{tcolorbox}[colback=infoblue!5!white,colframe=infoblue,title={\textbf{Este es el tutorial final de la hackathon}}]
Las 7 fases est\'an completas. El proyecto est\'a listo para el demo y el env\'io.
\end{tcolorbox}

\tableofcontents
\newpage

% ============================================================
\section{Estado del proyecto (al inicio de Fase 7)}
% ====================================

\subsection{Lo que ten\'iamos}

\begin{center}
\begin{tabular}{ll}
\toprule
\textbf{Fase} & \textbf{Estado} \\
\midrule
0. OCI setup & [OK] \\
1. Backend Core & [OK] \\
2. LLM integration & [OK] \\
3. Frontend deploy & [OK] \\
4. File uploads & [OK] \\
5. RAG con OCI & [OK] \\
6. Testing (50/50) & [OK] \\
7. Deploy final & \textbf{en progreso} \\
\bottomrule
\end{tabular}
\end{center}

\subsection{Lo que faltaba para ``production ready''}

\begin{itemize}
  \item \err \textbf{Bugs encontrados en testing} (4) $\rightarrow$ aplicar fix al c\'odigo deployado
  \item \err \textbf{Verificar E2E} que todo funciona despu\'es de los fixes
  \item \err \textbf{README} completo para el repositorio
  \item \err \textbf{Documentar} el estado final del stack
\end{itemize}

% ============================================================
\section{Paso 1: Re-deploy con los fixes}
% ====================================

\subsection{¿Por qu\'e re-deploy?}

En Fase 6 encontramos 4 bugs gracias a los tests:
\begin{enumerate}
  \item \texttt{query\_by\_tenant} usaba \texttt{eq} en vez de \texttt{begins\_with}
  \item \texttt{ScanIndexForward} no soportado en \texttt{scan}
  \item SyntaxError en \texttt{chat.py} (\texttt{\textbackslash{}\~{}})
  \item Imports duplicados (clases diferentes)
\end{enumerate}

Aunque el bug \#4 era solo de tests, los otros 3 afectaban el c\'odigo de producci\'on. Re-deployamos.

\subsection{Comando de re-deploy}

\begin{lstlisting}[language=bash,caption={Re-deploy con todos los fixes}]
cd /home/lFunknown/sentinel-academia/infra

# 1. Limpiar build anterior (si hay problemas de permisos)
rm -rf .aws-sam-new

# 2. Build con el c\'odigo actualizado
sam build --build-dir .aws-sam-new
# Output esperado:
#   Built Artifacts  : .aws-sam-new
#   Built Template   : .aws-sam-new/template.yaml

# 3. Deploy
sam deploy \
  --stack-name sentinel-academia-dev \
  --template-file .aws-sam-new/template.yaml \
  --s3-bucket sentinel-sam-artifacts-1781951709 \
  --capabilities CAPABILITY_IAM \
  --no-disable-rollback \
  --no-confirm-changeset
# Output esperado:
#   Successfully created/updated stack - sentinel-academia-dev
\end{lstlisting}

% ============================================================
\section{Paso 2: Test E2E completo (5 pasos)}
% ====================================

\subsection{El test que prueba TODO}

\begin{lstlisting}[language=bash,caption={Test E2E final completo}]
API="https://om3eiczjr9.execute-api.us-east-1.amazonaws.com/dev"

# Paso 1: Crear queja
echo "=== 1. POST /api/quejas ==="
RESP=$(curl -s -X POST "$API/api/quejas" \
  -H "Content-Type: application/json" \
  -H "X-Tenant-ID: demo-utpl" \
  -H "X-Correlation-ID: final-test-001" \
  -d '{"titulo":"Test final del proyecto",
       "descripcion":"Esta queja prueba el flujo completo",
       "categoriaDeclarada":"INFRAESTRUCTURA",
       "anonima":false}')
echo "$RESP" | python3 -m json.tool
QID=$(echo "$RESP" | python3 -c "import sys,json;print(json.load(sys.stdin)['quejaId'])")

# Paso 2: Pedir URL para subir archivo
echo "=== 2. POST /api/quejas/{id}/upload-url ==="
URL_RESP=$(curl -s -X POST "$API/api/quejas/$QID/upload-url" \
  -H "Content-Type: application/json" \
  -H "X-Tenant-ID: demo-utpl" \
  -d '{"filename":"evidencia.pdf","contentType":"application/pdf","fileSize":2048}')
UPLOAD_URL=$(echo "$URL_RESP" | python3 -c "import sys,json;print(json.load(sys.stdin)['uploadUrl'])")
echo "  Upload URL: ${UPLOAD_URL:0:80}..."

# Paso 3: Subir archivo a S3
echo "=== 3. PUT a S3 (presigned URL) ==="
echo "test" > /tmp/evidencia.pdf
curl -s -X PUT "$UPLOAD_URL" \
  -H "Content-Type: application/pdf" \
  --data-binary @/tmp/evidencia.pdf \
  -w "  HTTP %{http_code}\n" -o /dev/null

# Paso 4: Hacer pregunta al chat (RAG)
echo "=== 4. POST /api/chat ==="
sleep 3
curl -s -X POST "$API/api/chat" \
  -H "Content-Type: application/json" \
  -H "X-Tenant-ID: demo-utpl" \
  -d '{"question":"Cuales son las categorias de quejas?"}' | python3 -m json.tool | head -20

# Paso 5: Verificar queja completa
echo "=== 5. GET /api/quejas/{id} (final) ==="
sleep 3
curl -s "$API/api/quejas/$QID" -H "X-Tenant-ID: demo-utpl" | python3 -m json.tool | head -20
\end{lstlisting}

\subsection{Output esperado (cada paso)}

\subsubsection{Paso 1: POST /api/quejas}

\begin{lstlisting}[language=json]
{
  "quejaId": "e096f2e6-4a73-41a6-8fca-d9c0a180589f",
  "status": "PENDIENTE",
  "correlationId": "final-test-001",
  "createdAt": "2026-06-20T15:51:22.849474+00:00",
  "estimatedAnalysisTime": 30
}
\end{lstlisting}

\subsubsection{Paso 3: PUT a S3}

\begin{lstlisting}
HTTP 200
\end{lstlisting}

\subsubsection{Paso 4: POST /api/chat}

\begin{lstlisting}[language=json]
{
  "answer": "Basado en el reglamento, las categorias son...",
  "sources": [
    {"id": "reglamento-cap1", "title": "Capitulo 1: ...", "score": 0.5},
    {"id": "reglamento-cap4", "title": "Capitulo 4: ...", "score": 0.3},
    {"id": "reglamento-faq-acoso", "title": "FAQ: ...", "score": 0.2}
  ],
  "modeloUsado": "cohere.command-latest",
  "tokensUsados": 146
}
\end{lstlisting}

\subsubsection{Paso 5: GET /api/quejas/{id}}

\begin{lstlisting}[language=json]
{
  "quejaId": "e096f2e6-...",
  "status": "ANALIZADA",
  "adjuntos": [
    "https://sentinel-files-dev-...s3.us-east-1.amazonaws.com/..."
  ],
  "analysis": {
    "categoria": "INFRAESTRUCTURA",
    "criticidad": "MEDIA",
    "prioridad": 5,
    "tokensConsumidos": 25
  }
}
\end{lstlisting}

\subsection{¿Qu\'e validamos?}

\begin{itemize}
  \item \ok \textbf{Crear queja}: 202 + UUID + status PENDIENTE
  \item \ok \textbf{Pedir presigned URL}: retorna URL v\'alida con key multi-tenant
  \item \ok \textbf{Upload a S3}: 200 (sin auth, gracias a presigned URL)
  \item \ok \textbf{Chat (RAG)}: respuesta con sources y scores
  \item \ok \textbf{Status workflow}: PENDIENTE $\rightarrow$ ANALIZADA
  \item \ok \textbf{Adjuntos}: archivo subido v\'ia S3 aparece en la queja
  \item \ok \textbf{Analysis}: el mock LLM clasific\'o correctamente
\end{itemize}

% ============================================================
\section{Paso 3: Listar todos los recursos del stack}
% ====================================

\subsection{Comando para listar}

\begin{lstlisting}[language=bash,caption={Listar recursos del stack}]
aws cloudformation describe-stack-resources \
  --stack-name sentinel-academia-dev \
  --query 'StackResources[].{Type:ResourceType, Name:LogicalResourceId, Id:PhysicalResourceId}' \
  --output table
\end{lstlisting}

\subsection{Output (resumen)}

\begin{center}
\begin{tabular}{lll}
\toprule
\textbf{Tipo} & \textbf{Nombre} & \textbf{ID/ARN} \\
\midrule
AWS::Lambda::Function & HealthCheckFunction & sentinel-health-dev \\
AWS::Lambda::Function & CreateQuejaFunction & sentinel-create-queja-dev \\
AWS::Lambda::Function & GetQuejaFunction & sentinel-get-queja-dev \\
AWS::Lambda::Function & ListQuejasFunction & sentinel-list-quejas-dev \\
AWS::Lambda::Function & AnalyzeQuejaFunction & sentinel-analyze-queja-dev \\
AWS::Lambda::Function & UploadURLFunction & sentinel-upload-url-dev \\
AWS::Lambda::Function & FileProcessorFunction & sentinel-file-processor-dev \\
AWS::Lambda::Function & ChatFunction & sentinel-chat-dev \\
\midrule
AWS::DynamoDB::Table & QuejasTable & sentinel-quejas-dev \\
AWS::SQS::Queue & AnalysisQueue & sentinel-analysis-dev \\
AWS::SQS::Queue & AnalysisDeadLetterQueue & sentinel-analysis-dlq-dev \\
AWS::S3::Bucket & FilesBucket & sentinel-files-dev-227165337884 \\
AWS::ApiGateway::RestApi & SentinelApi & om3eiczjr9 \\
\bottomrule
\end{tabular}
\end{center}

\subsection{Outputs del stack}

\begin{lstlisting}[language=bash,caption={Obtener outputs}]
aws cloudformation describe-stacks \
  --stack-name sentinel-academia-dev \
  --query 'Stacks[0].Outputs' \
  --output table
\end{lstlisting}

\begin{lstlisting}[caption={Output esperado}]
--------------------------------------------------------------------------------------
|  OutputKey          |  OutputValue
--------------------------------------------------------------------------------------
|  ApiUrl             |  https://om3eiczjr9.execute-api.us-east-1.amazonaws.com/dev
|  QuejasTableName    |  sentinel-quejas-dev
|  AnalysisQueueUrl   |  https://sqs.us-east-1.amazonaws.com/.../sentinel-analysis-dev
|  FilesBucketName    |  sentinel-files-dev-227165337884
|  FrontendUrl        |  http://sentinel-frontend-dev-...s3-website-us-east-1.amazonaws.com
--------------------------------------------------------------------------------------
\end{lstlisting}

% ============================================================
\section{Paso 4: Verificar tests}
% ====================================

\subsection{Correr pytest}

\begin{lstlisting}[language=bash,caption={Verificar 50/50 tests pasando}]
cd /home/lFunknown/sentinel-academia
export PYTHONPATH=src:$PYTHONPATH
python3 -m pytest src/tests/ -v
\end{lstlisting}

\subsection{Output esperado}

\begin{lstlisting}[caption={}]
src/tests/unit/test_shared.py::TestConfig::test_default_values PASSED
src/tests/unit/test_shared.py::TestConfig::test_dynamodb_table_resolved PASSED
src/tests/unit/test_shared.py::TestCorrelationId::test_with_correlation_id PASSED
... (47 more)
src/tests/integration/test_handlers.py::TestChat::test_chat_requires_tenant PASSED

============================== 50 passed in 2.09s ==============================
\end{lstlisting}

% ============================================================
\section{Paso 5: README.md para el repositorio}
% ====================================

\subsection{El README completo}

El \texttt{README.md} en la ra\'iz del proyecto incluye:

\begin{itemize}
  \item \textbf{Arquitectura} con diagrama ASCII
  \item \textbf{Stack tecnol\'ogico} completo (Python, AWS, OCI, Vue)
  \item \textbf{Estructura del proyecto} (todos los archivos)
  \item \textbf{Setup local} paso a paso
  \item \textbf{Deploy en AWS} con comandos
  \item \textbf{Configuraci\'on de OCI} con Python SDK
  \item \textbf{Endpoints de la API} con ejemplos curl
  \item \textbf{Tests} c\'omo correrlos
  \item \textbf{Troubleshooting} de errores comunes
  \item \textbf{Stack final} desplegado (todos los recursos)
\end{itemize}

\subsection{Fragmentos clave del README}

\begin{lstlisting}[caption={README.md — Instalacion}]
## Setup local

### Requisitos
- Python 3.12+
- Node.js 18+
- AWS CLI 2.x
- SAM CLI 1.162+
- mise 2026+
- Tectonic 0.16+ (para compilar PDFs)

### 1. Clonar
git clone https://github.com/your-org/sentinel-academia.git
cd sentinel-academia

### 2. Python 3.12
mise use --global python@3.12

### 3. Deps Python
pip install pydantic pydantic-settings email-validator boto3 \
  aws-lambda-powertools[tracer,logger,metrics] orjson \
  pytest pytest-cov pytest-mock 'moto[dynamodb,s3,sqs]' \
  boto3-stubs[dynamodb,s3,sqs]

### 4. Deps Node
cd web && npm install && cd ..

### 5. Credenciales AWS Academy
export AWS_ACCESS_KEY_ID="ASIAT..."
export AWS_SECRET_ACCESS_KEY="..."
export AWS_SESSION_TOKEN="IQoJb3JpZ2luX2Vj..."
export AWS_DEFAULT_REGION="us-east-1"
export AWS_REGION="us-east-1"
aws sts get-caller-identity
\end{lstlisting}

\begin{lstlisting}[caption={README.md — Deploy}]
## Deploy en AWS

### Build
cd infra
sam build --build-dir .aws-sam-new

### Deploy
sam deploy \
  --stack-name sentinel-academia-dev \
  --template-file .aws-sam-new/template.yaml \
  --s3-bucket YOUR-SAM-ARTIFACTS-BUCKET \
  --capabilities CAPABILITY_IAM \
  --no-disable-rollback \
  --no-confirm-changeset

### S3 notification (post-deploy)
BUCKET=$(aws cloudformation describe-stacks \
  --stack-name sentinel-academia-dev \
  --query 'Stacks[0].Outputs[?OutputKey==`FilesBucketName`].OutputValue' \
  --output text)

LAMBDA_ARN=$(aws lambda get-function \
  --function-name sentinel-file-processor-dev \
  --query 'Configuration.FunctionArn' --output text)

aws s3api put-bucket-notification-configuration \
  --bucket "$BUCKET" \
  --notification-configuration "{...}"
\end{lstlisting}

% ============================================================
\section{Cat\'alogo de errores y soluciones (Fase 7)}
% ====================================

\begin{tcolorbox}[colback=errred!5!white,colframe=errred,breakable,title={Error 1: Tests pasan pero deploy falla}]
\textbf{Cuando}: \texttt{pytest} OK pero el sistema en prod no. \\
\textbf{Causa}: tests usan moto (mock), prod usa AWS real. \\
\textbf{Soluci\'on}: hacer test E2E con curl al API real.
\end{tcolorbox}

\begin{tcolorbox}[colback=errred!5!white,colframe=errred,breakable,title={Error 2: CORS preflight fails}]
\textbf{Cuando}: el frontend no puede llamar al API. \\
\textbf{Causa}: API Gateway no tiene headers CORS. \\
\textbf{Soluci\'on}: verificar \texttt{Globals.Api.Cors} en el template.
\end{tcolorbox}

\begin{tcolorbox}[colback=errred!5!white,colframe=errred,breakable,title={Error 3: S3 notification no dispara}]
\textbf{Cuando}: subes archivo pero la queja no se actualiza. \\
\textbf{Causa}: notification config no aplicada. \\
\textbf{Soluci\'on}: \texttt{aws s3api get-bucket-notification-configuration}.
\end{tcolorbox}

\begin{tcolorbox}[colback=errred!5!white,colframe=errred,breakable,title={Error 4: Stack en ROLLBACK\_COMPLETE}]
\textbf{Cuando}: deploy falla y no se puede re-deployar. \\
\textbf{Soluci\'on}:
\begin{lstlisting}
aws cloudformation delete-stack --stack-name sentinel-academia-dev
aws cloudformation wait stack-delete-complete --stack-name sentinel-academia-dev
sam deploy ...
\end{lstlisting}
\end{tcolorbox}

% ============================================================
\section{Lo que aprendimos (resumen ejecutivo)}
% ====================================

\begin{tcolorbox}[colback=okgreen!5!white,colframe=okgreen,title={\textbf{Checklist Fase 7}}]
\begin{enumerate}
  \item \textbf{Re-deploy} despu\'es de tests para aplicar fixes
  \item \textbf{Test E2E} con curl (5 pasos secuenciales)
  \item \textbf{Verificar} todos los outputs del stack
  \item \textbf{pytest} sigue verde despu\'es del re-deploy
  \item \textbf{README.md} completo para el repositorio
  \item \textbf{Documentar} el estado final (todos los recursos)
  \item \textbf{Tests pasan} (50/50)
  \item \textbf{7 fases} completas
  \item \textbf{Hackathon} lista para el demo
  \item \textbf{Siguiente}: el usuario graba el video demo
\end{enumerate}
\end{tcolorbox}

% ============================================================
\section{Siguiente paso: tu video demo}
% ====================================

El usuario grabar\'a el video demo. Aqu\'i hay un gui\'on sugerido:

\begin{enumerate}
  \item \textbf{Intro (15s)}: ``Soy [nombre], esto es Sentinel AcademIA, una plataforma
        de quejas universitarias con IA. Hackathon 2026.''
  \item \textbf{Arquitectura (30s)}: mostrar el diagrama en el README.
        Explicar AWS + OCI.
  \item \textbf{Demo en vivo (90s)}:
    \begin{itemize}
      \item Abrir el frontend (\texttt{http://sentinel-frontend-dev-...s3-website...})
      \item Crear una queja: ``Problema con el aula 205''
      \item Esperar 5 segundos
      \item Mostrar el an\'alisis: categor\'ia INFRAESTRUCTURA, criticidad MEDIA
      \item Hacer una pregunta al chat: ``\`Qu\'e hago si sufro acoso?''
      \item Mostrar la respuesta con sources del reglamento
    \end{itemize}
  \item \textbf{Backend (30s)}: mostrar la consola de AWS:
    \begin{itemize}
      \item CloudWatch logs de la Lambda
      \item DynamoDB con la queja creada
      \item S3 con el archivo de evidencia
    \end{itemize}
  \item \textbf{Tests (15s)}: mostrar \texttt{pytest} corriendo 50 tests en 2 segundos.
  \item \textbf{Cierre (15s)}: ``50 tests, 7 Lambdas, multi-tenant, RAG, y todo deployado
        en AWS Academy. Gracias.''
\end{enumerate}

\vspace{2em}
\begin{center}
\textit{Fin del tutorial. \`Exito en tu demo!}
\end{center}

\end{document}
